\section{Moment relaxations and prefix separation}\label{sec:setup}

We first describe the finite collection of statistics used by the decoder
and the convex set against which they will separate the true input.
Throughout, $[n]=\{1,\ldots,n\}$ and $X_1,\ldots,X_n$ are formal Boolean
variables. Products and degrees are understood after the reductions
$X_i^2=X_i$.

\subsection{Observable trace moments}
For a nonempty rank set $I=\{\ell_1<\cdots<\ell_k\}\subseteq[n]$,
the corresponding trace observable is the product of the bits at those
ranks, with value zero if the trace has length less than $\ell_k$.
Let $T_I(X)$ be its channel expectation, regarded as a polynomial in
the input bits. More explicitly,
\[
 T_I(X)=\sum_{i_1<\cdots<i_k}
    c_{I;i_1,\ldots,i_k}\,X_{i_1}\cdots X_{i_k},
\]
where $c_{I;i_1,\ldots,i_k}$ is the probability that the selected input
positions appear at ranks $\ell_1,\ldots,\ell_k$ in a deletion trace.
These events are disjoint as the input tuple varies. Consequently the
coefficients are nonnegative, their sum is at most one, and the degree
is at most $k$. They are rational when $p$ is rational. The algorithm
will estimate $T_I(x)$ by averaging the observable over independent
traces; it never observes the input positions that determine the
coefficients.

\subsection{Positive moment arrays}

The following Boolean moment relaxation is a specialization of the
semidefinite moment and sum-of-squares frameworks of
Lasserre~\cite[Section~2]{Lasserre2001} and
Parrilo~\cite[Sections~4.2--4.3]{Parrilo2000}. We use only its defining
positivity condition and elementary consequences, proved below.
Fix a proposed prefix $u\in\{0,1\}^d$ and an integer $D\ge1$.
On the remaining positions $U=\{d+1,\ldots,n\}$, use sign variables
$Y_i=2X_i-1$, so $Y_i^2=1$. A moment array consists of real numbers
$m_J$ for $J\subseteq U$ with $|J|\le2D$, including $m_\varnothing=1$,
subject to
\begin{equation}\label{eq:moment-matrix}
 \mathcal M_D(m)=
 \bigl(m_{I\triangle J}\bigr)_{
       I,J\subseteq U,\ |I|,|J|\le D}\succeq0.
\end{equation}
Only existing subsets are indexed; thus the definition also applies
when $D>|U|$ or $U$ is empty.

The array defines a functional $\EE$ on polynomials of degree at most
$2D$ by substituting $X_i=u_i$ for
$i\le d$, expanding the other bits as $(1+Y_i)/2$, reducing $Y_i^2=1$,
and setting $\EE[\prod_{i\in J}Y_i]=m_J$. We extend this functional
complex linearly. We call it a \emph{positive moment functional of order
$2D$ fixing $u$}. This terminology asserts \eqref{eq:moment-matrix},
not the existence of a probability distribution realizing all its moments.

\begin{lemma}[Elementary moment bounds]\label{lem:moment-bounds}
For every array satisfying \eqref{eq:moment-matrix}, one has
$|m_J|\le1$ for $|J|\le2D$. The value under $\EE$ of any product of at
most $2D$ distinct input bits has magnitude at most one. On polynomials
of degree at most $D$, the form
\[
 \langle P,R\rangle_{\EE}=\EE[P\overline R]
\]
is a positive semidefinite Hermitian form, linear in its first argument.
It satisfies Cauchy--Schwarz and becomes an inner product after
quotienting by its zero-norm subspace.
\end{lemma}

\begin{proof}
The empty moment equals one. Every nonempty $J$ with $|J|\le2D$ can be partitioned as $J=I\sqcup K$ with
$|I|,|K|\le D$. The $I,K$ principal submatrix in
\eqref{eq:moment-matrix} has diagonal entries one and off-diagonal
entry $m_J$, giving $|m_J|\le1$. Expanding a product of unfixed bits
as $2^{-t}\prod(1+Y_i)$ is a convex combination of these sign monomials;
fixed bits either remove factors or make the product zero. This proves
the second assertion. The final assertion is the Gram-form
interpretation of the real symmetric positive semidefinite matrix
\eqref{eq:moment-matrix}, extended to complex coefficient vectors.
Cauchy--Schwarz follows by applying positivity to $P+zR$ for $z\in\CC$;
it also shows that every zero-norm vector is orthogonal to every vector,
so the quotient is well defined.
\end{proof}

Fix a true input $x\in\{0,1\}^n$ and a prefix $u$ of length $d$ with
\[
 u_i=x_i\quad(i<d),\qquad u_d=1-x_d.
\]
For a positive moment functional fixing $u$, write
\[
 \Delta[P]=\EE[P]-P(x).
\]
Then $\Delta[X_d]=\pm1$, and $\Delta[P]=0$ whenever $P$ uses only
positions strictly before $d$. If $P$ is a linear combination of bit
monomials of degree at most $2D$, Lemma~\ref{lem:moment-bounds} gives
\begin{equation}\label{eq:delta-coefficients}
 |\Delta[P]|\le2\sum_A |c_A|
 \quad\text{when }P=\sum_A c_A\prod_{i\in A}X_i.
\end{equation}

For $p<1$, introduce
\begin{equation}\label{eq:aux-parameters}
 q_* =\max(q,1/n),\qquad
 Q=1+\log(1/q_*),\qquad
 Y=1+\frac{\log n}{Q}.
\end{equation}
When $q\ge1/n$, $Y=\mathsf X(n,p)$. When $q<1/n$,
$1\le\mathsf X(n,p)\le Y<2$. Hence
\[
 \mathsf X(n,p)\le Y\le2\mathsf X(n,p),\qquad
 QY=Q+\log n\le C\log n
\]
for an absolute constant $C$, since $n\ge2$.

\begin{theorem}[Separation from a wrong prefix]\label{thm:separation}
There exist absolute constants $C_{\mathrm{sep}}\ge1$ and an integer
$H\ge1$ with the following property. For $n\ge2$ and $0<p<1$, define
\[
 D_*=4^{\lceil\log_2 Y\rceil+H}\le4^{H+1}Y^2,
\]
using \eqref{eq:aux-parameters}. For every $x\in\{0,1\}^n$, every
$d\in[n]$, and every positive moment functional of order $2D$ fixing
$x_1,\ldots,x_{d-1},1-x_d$, where $D\ge D_*$, some nonempty
$I\subseteq[n]$ with $|I|\le D_*$ satisfies
\[
 |\EE[T_I]-T_I(x)|>
       \exp\bigl(-C_{\mathrm{sep}}B(n,p)\bigr).
\]
\end{theorem}

The theorem concerns every feasible array, including arrays not
represented by distributions. Its proof occupies the next two sections.
Once it is available, empirical trace moments rule out every wrong
extension of a correct prefix. Section~\ref{sec:decoding} shows how to
find the feasible extension with an explicit bounded computation.

\subsection{Exponentially weighted tuple statistics}
For the analytic proof, fix $x,d,\EE$ as above. For $1\le k\le2D$,
an \emph{order-$k$ feature}
is a sequence of polynomials
\begin{equation}\label{eq:feature}
 f_j=\sum_{i_1<\cdots<i_k=j}
       X_{i_1}\cdots X_{i_k}
       \prod_{h=1}^{k-1}z_h^{i_{h+1}-i_h},
 \qquad j\in[n],
\end{equation}
where $z_h=\exp(-a_h+\ii\theta_h)$, $a_h>0$, and $\theta_h\in\RR$.
For $k=1$ this means $f_j=X_j$. Put
\[
 H_f=\prod_{h=1}^{k-1}(1+1/a_h).
\]
Summing geometric series in the positive gaps shows that $H_f\ge1$
bounds the sum of the absolute coefficients of each $f_j$, since
$(1-e^{-a})^{-1}\le1+1/a$ for $a>0$.
For $s>0$ and $\omega\in\RR$, define its signal by
\begin{equation}\label{eq:signal}
 \mathcal S_f(s,\omega)=
 \sum_{j=1}^n\Delta[f_j]e^{(-s+\ii\omega)(j-d)}.
\end{equation}
Every monomial in $f_j$ uses positions at most $j$, so
$\Delta[f_j]=0$ for $j<d$. Thus, despite the centered exponent,
only nonnegative distances contribute to \eqref{eq:signal}.
The purpose of the multiscale argument is to preserve a detectable
signal while reducing $s$ and $|\omega|$, at the cost of increasing $k$.
